\documentclass[acmsmall,nonacm]{acmart}

\settopmatter{printacmref=false, printccs=false, printfolios=true}
\setcopyright{none}
\renewcommand\footnotetextcopyrightpermission[1]{}

\citestyle{acmauthoryear}

\usepackage{booktabs}
\usepackage{listings}
\usepackage{xcolor}

\lstset{
  basicstyle=\ttfamily\footnotesize,
  columns=fullflexible,
  keepspaces=true,
  breaklines=true,
  frame=single,
  framerule=0.2pt,
  xleftmargin=0.5em,
  xrightmargin=0.5em
}

\newcommand{\Rocq}{Rocq}
\newcommand{\Coq}{Coq}
\newcommand{\CertiGC}{CertiGC}
\newcommand{\CertiGraph}{CertiGraph}
\newcommand{\VST}{VST}
\newcommand{\todo}[1]{\textcolor{red}{\textbf{TODO:} #1}}

\title{AI-Assisted Completion of \CertiGC{} Proofs: An Experience Report}

\author{Shengyi Wang}
\affiliation{%
  \institution{Shanghai Qi Zhi Institute}
  \city{Shanghai}
  \country{China}
}
\email{wangshengyi@sqz.ac.cn}

\begin{abstract}
This experience report describes the use of OpenAI Codex to complete and
stabilize a substantial \Rocq{}/\Coq{} proof development for \CertiGC{}, the
verified generational garbage collector in the \CertiGraph{} project.  The
development extends the collector from an effectively immutable setting to a
mutable one by adding remembered-set forwarding to the collection path and
re-establishing the top-level graph-isomorphism correctness theorem.  The
central technical issue was not low-level proof scripting alone: mutable updates
invalidate the old global \texttt{no\_backward\_edge} assumption, so the proof
had to be reorganized around a remembered-set invariant stating that every
backward edge is recorded.

The case differs from recent reports of LLMs generating new compiler proofs
from a close template.  Here Codex was used to finish a long-running, partially
developed verification in a mature and brittle \VST{}/\CertiGraph{} codebase.
The machine checker remained the arbiter of correctness, while the human role
shifted toward selecting invariants, constraining specification changes,
reviewing theorem statements, and deciding when proof cleanup was justified.
We report the workflow, the resulting artifact, and lessons about where an
agentic assistant was useful and where human proof-engineering judgment remained
essential.
\end{abstract}

\begin{document}

\maketitle

\section{Introduction}

Large language models are now routinely used as programming assistants, but
their usefulness for formal verification remains more delicate than for ordinary
software development.  An assistant can generate a large volume of code quickly,
yet the risk is not merely that the code fails to compile.  In a proof
assistant, the most damaging failure mode is often a silently weakened theorem,
an inappropriate new precondition, or a specification change that makes the
proved result no longer say what the development needs.

This paper reports on a concrete proof-engineering experience in that setting:
using OpenAI Codex to help complete the mutable-garbage-collector branch of
\CertiGC{} \citep{openai2025codex,openai2026codexapp}.  \CertiGC{} is part of the \CertiGraph{} development, which uses
\VST{} and \Coq{}/\Rocq{} to verify C programs manipulating heap-represented
graphs \citep{wang2019certigraph,appel2014programlogics}.  The collector is a
generational copying collector.  The mutable extension adds a write barrier and
remembered sets so that the collector remains correct when older objects can be
mutated to point to younger ones.

The work began as a conventional human proof effort.  The early branch already
contained substantial design and verification work: remembered-set data
structures, a verified \texttt{forward\_remset}, updates to the mutator-side
\texttt{mutable\_update} specification, and a refactoring of thread information
to carry remembered-set state.  The later phase used Codex as an agentic coding
assistant: it read the codebase, edited proof files, ran \texttt{make}, diagnosed
errors, proposed helper lemmas, and performed cleanup.  The goal was not to
obtain a plausible proof script, but to get the repository into a state where
the existing verification target compiled and the final theorem had the intended
shape.

The experience was inspired by Paraskevopoulou's recent report on using Claude
Code to mechanize an ANF transformation correctness proof for CertiCoq
\citep{paraskevopoulou2026machinegenerated}.  That paper showed that an
agentic assistant can adapt an existing compiler-proof technique to a new
transformation and produce a machine-checked proof in days.  The present case
studies a different point in the design space.  Codex did not start from a clean
template and a fresh theorem.  Instead, it worked inside an existing
systems-verification development where the proof broke because the collector's
semantic invariant changed.

This report makes three contributions:

\begin{itemize}
\item It describes the proof task needed to make mutable \CertiGC{} compatible
with remembered sets while preserving the roots-level graph-isomorphism theorem.
\item It reports the Codex-assisted workflow used to finish the proof, including
how human guidance constrained specification changes and invariant design.
\item It compares this experience with a recent LLM-assisted compiler-proof
experience report and identifies lessons specific to maintaining large
\VST{}/\CertiGraph{} proofs.
\end{itemize}

\section{Background: \CertiGC{}, \VST{}, and Remembered Sets}

\CertiGraph{} provides a framework for verifying graph-manipulating C programs:
abstract graphs are represented in \Coq{}, while separation-logic predicates
connect those graphs to concrete heap layouts \citep{wang2019certigraph}.
\CertiGC{} is the generational garbage collector in that development.  Its proof
uses \VST{}'s Verifiable C program logic, whose soundness is connected to
CompCert's C semantics through the Verified Software Toolchain
\citep{appel2014programlogics,leroy2009compcert}.

Generational collection relies on the observation that most objects die young.
The collector copies live objects from a younger ``from'' generation to an older
``to'' generation, forwarding roots and then scanning copied objects until all
reachable objects have been copied.  In the original proof structure, the heap
was effectively immutable from the collector's perspective: the proof could rely
on the absence of old-to-young edges, usually stated as a
\texttt{no\_backward\_edge}-style invariant.

Mutable updates invalidate that invariant.  A program may update a field in an
older object to point to a younger object.  A generational collector can still be
correct if such edges are recorded by a write barrier.  During collection, the
remembered set is treated as an additional source of effective roots: the
collector must forward objects reachable from recorded old-to-young edges before
ordinary root forwarding and scanning are enough.

In the C code in this branch, the generation collection path therefore starts by
forwarding remembered-set entries:

\begin{lstlisting}[language=C]
value *p = to->next;
forward_remset(from, to, &to->next);
forward_roots(from->start, from->limit, &to->next, fr);
do_scan(from->start, from->limit, p, &to->next);
\end{lstlisting}

At the specification level, \texttt{garbage\_collect\_condition} remains a
heap/graph condition:

\begin{lstlisting}
Definition garbage_collect_condition (g: LGraph) (h : part_heap) : Prop :=
  graph_unmarked g /\ no_dangling_dst g /\ ti_size_spec h.
\end{lstlisting}

The remembered-set condition is kept separate:

\begin{lstlisting}
Definition no_unrecorded_backward_edge (g: LGraph) (rh: remset_heap): Prop :=
  forall e,
    graph_has_e g e ->
    (egeneration e > vgeneration (dst g e))%nat ->
    In (RemSetInterior
          (InteriorVertexPos (fst e) (Z.of_nat (snd e))))
       (nth_remset_space rh (vgeneration (dst g e))).
\end{lstlisting}

This separation became an important proof-engineering constraint.  Remembered
sets are internal collector state, and exposing them in the final user-facing
mathematical theorem would make the theorem less useful.  The proof was
therefore allowed to use remembered sets internally, but the end result still had
to be a roots-level graph-isomorphism statement.

\section{The Proof Task}

The final theorem in the branch is:

\begin{lstlisting}
Theorem garbage_collect_isomorphism:
  forall roots1 roots2 g1 h1 rh1 rmst1 g2 h2 rh2 rmst2,
    graph_unmarked g1 -> no_unrecorded_backward_edge g1 rh1 ->
    no_dangling_dst g1 ->
    roots_graph_compatible roots1 g1 ->
    sound_gc_graph g1 ->
    remset_graph_state g1 O rmst1 rh1 ->
    remset_heap_covers_graph g1 rh1 ->
    garbage_collect_relation roots1 roots2
      g1 h1 rh1 rmst1 g2 h2 rh2 rmst2 ->
    gc_graph_iso g1 roots1 g2 roots2.
\end{lstlisting}

The statement reveals the main technical compromise.  The theorem assumes the
remembered-set invariant and enough compatibility between remembered-set state
and the graph, but the conclusion is still the original graph-isomorphism
property.  The remembered set is a way to prove collection correct in the
presence of mutation; it is not part of the observable correctness result.

The hard part was to bridge remembered-set forwarding, root forwarding,
scanning, and reset into the graph-isomorphism proof.  Forwarding
remembered-set entries temporarily
turns recorded old-to-young edges into effective roots.  After the collection of
one generation, the remembered-set component for that generation is reset.  The
proof must show that no unrecorded backward edge is introduced by this sequence.
One key preservation lemma has the following role:

\begin{lstlisting}
do_generation_relation_no_unrecorded_backward_edge_reset
\end{lstlisting}

This lemma connects the full one-generation collection relation, including
remembered-set forwarding and reset, back to the invariant needed by the outer
garbage-collection loop.

The task was therefore not just to patch broken proof scripts.  It required a
change in proof architecture:

\begin{itemize}
\item remove the invalid global \texttt{no\_backward\_edge} route;
\item introduce and preserve \texttt{no\_unrecorded\_backward\_edge};
\item make remembered-set entries act as temporary roots during generation
copying;
\item keep the top-level theorem framed as graph isomorphism over ordinary
roots;
\item repair the \VST{} verification files so that the C code and mathematical
relation agree.
\end{itemize}

\section{Workflow with Codex}

The Codex-assisted workflow was closer to supervised proof maintenance than to
one-shot proof generation.  The assistant was given access to the repository and
was allowed to edit files and run targeted builds.  The human supplied the proof
goal, design constraints, and course corrections.

\paragraph{Development loop.}
Most iterations followed a simple pattern.  Codex inspected the current proof
failure, searched for nearby lemmas, edited a focused region, and ran a targeted
build such as:

\begin{lstlisting}
make CertiGC/GCGraph.vo
make CertiGC/gc_correct.vo
make CertiGC/verif_garbage_collect.vo
\end{lstlisting}

Full builds with \texttt{make -j8} were used at integration points.  The
machine checker was the primary feedback mechanism; Codex's natural-language
explanations were useful mainly when the failure indicated a mismatch between
the intended invariant and the proof state.

\paragraph{Human constraints.}
The human guidance was deliberately conservative.  New or changed lemmas were
acceptable, but new or changed \texttt{Definition}s required discussion.
Preconditions were not to be added to \VST{} specs unless there was concrete
evidence that the source program required them.  In particular,
\texttt{no\_unrecorded\_backward\_edge} was kept as a separate
\texttt{garbage\_collect\_spec} precondition rather than being silently folded
into \texttt{garbage\_collect\_condition}.  This mattered because an LLM can make
progress by strengthening assumptions in ways that compile but weaken the
scientific result.

\paragraph{What Codex was good at.}
Codex was effective at local proof repair and repetitive invariant propagation.
Once the remembered-set invariant and bridge shape were clear, many subgoals
amounted to proving that a property was preserved by a fold over
\texttt{forward\_remset\_item}, by heap updates, or by adding/resetting one
generation.  These tasks are tedious for a human but well suited to an agent
that can search the development, try variants, and run the checker repeatedly.

Codex was also useful for cleanup after the theorem was restored.  The branch
contains several later commits whose purpose was not to change the theorem but
to remove stale proof leftovers, factor repeated preservation arguments, update
deprecated names, and simplify helper lemmas.  In a large proof development,
this maintenance is not cosmetic: stale helper lemmas obscure which invariants
are actually needed and make future repair harder.

\paragraph{Where human judgment remained central.}
The critical invariant choice was human-level proof engineering.  The old proof
failed because its semantic assumption was no longer true, not because a tactic
was missing.  Replacing \texttt{no\_backward\_edge} with a remembered-set
condition required understanding the collector algorithm, the write barrier, and
the desired external theorem.  Codex could help explore consequences and fill in
preservation lemmas, but the proof had to be guided away from easy but
undesirable exits, such as adding stronger heap/spatial preconditions to pure
specifications or exposing internal remembered-set details in the final
statement.

\todo{Add a short curated set of actual prompts from the Codex transcript, if
available and safe to quote.  The most useful examples would be: the initial
task framing, a prompt forbidding unnecessary preconditions, a prompt correcting
a mistaken proof direction, and a prompt asking for cleanup after the theorem
compiled.}

\section{Results}

The branch reached a compiling state and was merged into the upstream
\texttt{live} branch as pull request 30 with the merge message ``Replace invalid
\texttt{no\_backward\_edge} invariant for mutable GC'' on May 3, 2026.

Table~\ref{tab:repo-evidence} summarizes repository evidence extracted from the
git history.  These numbers should be read as artifact statistics, not as a
precise measure of human or AI labor.  The git author field does not distinguish
manual edits from Codex-assisted edits.

\begin{table}[t]
\centering
\caption{Repository evidence for the mutable-GC proof effort.}
\label{tab:repo-evidence}
\begin{tabular}{lll}
\toprule
Interval & Evidence & Notes \\
\midrule
Nov. 2024--Aug. 2025 & 16 CertiGC commits; +5292/-2086 & early groundwork \\
Mar. 26, 2026 & +30/-33 in one file & compatibility checkpoint \\
Apr. 22--May 3, 2026 & 51 CertiGC commits; +13340/-3557 & final sprint \\
Nov. 2024--May 3, 2026 & 74 CertiGC commits; +17953/-4967 & whole branch slice \\
\bottomrule
\end{tabular}
\end{table}

The final sprint touched 18 \CertiGC{} files.  The largest proof changes were in
\texttt{GCGraph.v}, which accumulated the remembered-set invariants and
preservation facts, and \texttt{gc\_correct.v}, which re-established the
graph-isomorphism bridge.  The verification-facing changes were concentrated in
\texttt{verif\_do\_generation.v}, \texttt{verif\_garbage\_collect.v}, and
\texttt{gc\_spec.v}.

In addition to repository statistics, I extracted interaction statistics from
the two Codex rollout transcripts corresponding to this work.  The parser
counts structured rollout items rather than streamed token deltas: user prompts
are \texttt{message} items with role \texttt{user}, excluding automatic
\texttt{<environment\_context>} messages; tool calls are \texttt{function\_call},
\texttt{custom\_tool\_call}, and \texttt{web\_search\_call} items; context
compactions are top-level \texttt{compacted} records.  Active time is computed by
summing gaps of at most 30 minutes between recorded rollout events, following
the same convention used by \citet{paraskevopoulou2026machinegenerated}.  Rocq
MCP calls are a subset of the function-tool-call count, not an additional
category added on top of total tool calls.  I report wall-clock time as the sum
of each active day's first-to-last transcript span; the longer calendar span is
included only to make clear that the transcript spans multiple calendar days.
The statistics are cut off at May 3, 2026, the end of the mutable-collector
development sprint; later transcript entries about writing this report are
excluded.

\begin{table}[t]
\centering
\caption{Transcript-derived Codex usage statistics.}
\label{tab:transcript-stats}
\begin{tabular}{lr}
\toprule
Metric & Value \\
\midrule
Human prompts & 415 \\
Context compactions & 88 \\
Total tool calls & 13,305 \\
Shell command calls & 8,077 \\
Shell session inputs & 1,368 \\
Rocq MCP calls & 2,096 \\
Proof-checking invocations & 2,131 \\
Shell build commands & 795 \\
Shell search commands & 2,018 \\
Shell read/inspection commands & 5,337 \\
Apply-patch edit calls & 1,731 \\
Git commit commands & 80 \\
Calendar-span hours & 265.9 \\
Wall-clock hours, per-day sum & 106.5 \\
Active hours & 59.3 \\
\bottomrule
\end{tabular}
\end{table}

Table~\ref{tab:per-day-stats} breaks the same categories down by day.

\begin{table}[t]
\centering
\caption{Per-day Codex development activity.}
\label{tab:per-day-stats}
\footnotesize
\setlength{\tabcolsep}{2pt}
\begin{tabular}{@{}l*{14}{r}@{}}
\toprule
Day & Prompts & Compact & Tools & Shell & Inputs & Rocq & Checks & Build & Search & Read & Edits & Git & Wall(h) & Active(h) \\
\midrule
Apr. 22 & 24 & 2 & 539 & 384 & 74 & 20 & 67 & 60 & 111 & 210 & 54 & 4 & 5.7 & 2.9 \\
Apr. 23 & 2 & 0 & 52 & 32 & 1 & 18 & 17 & 1 & 11 & 20 & 1 & 0 & 0.8 & 0.8 \\
Apr. 26 & 42 & 8 & 1,750 & 912 & 346 & 255 & 303 & 212 & 195 & 482 & 234 & 10 & 9.0 & 7.6 \\
Apr. 27 & 58 & 15 & 2,499 & 1,326 & 417 & 503 & 474 & 168 & 342 & 878 & 250 & 20 & 14.1 & 10.1 \\
Apr. 28 & 93 & 26 & 3,070 & 1,837 & 32 & 727 & 602 & 82 & 516 & 1,296 & 465 & 7 & 23.0 & 12.8 \\
Apr. 29 & 62 & 23 & 2,705 & 1,657 & 118 & 546 & 402 & 29 & 440 & 1,215 & 381 & 11 & 21.4 & 10.8 \\
Apr. 30 & 53 & 4 & 759 & 536 & 129 & 4 & 61 & 57 & 102 & 311 & 82 & 9 & 11.6 & 4.9 \\
May 1 & 43 & 6 & 1,031 & 686 & 211 & 0 & 110 & 110 & 153 & 447 & 134 & 5 & 10.8 & 4.9 \\
May 2 & 20 & 2 & 487 & 401 & 2 & 12 & 60 & 52 & 79 & 262 & 72 & 6 & 5.3 & 2.1 \\
May 3 & 18 & 2 & 413 & 306 & 38 & 11 & 35 & 24 & 69 & 216 & 58 & 8 & 4.8 & 2.3 \\
\midrule
Total & 415 & 88 & 13,305 & 8,077 & 1,368 & 2,096 & 2,131 & 795 & 2,018 & 5,337 & 1,731 & 80 & 106.5 & 59.3 \\
\bottomrule
\end{tabular}
\end{table}

The transcript-derived file activity aligns with the git history.  The most
heavily edited files were \texttt{gc\_correct.v} (836 patch calls),
\texttt{GCGraph.v} (384), \texttt{verif\_garbage\_collect.v} (226), and
\texttt{verif\_do\_generation.v} (131).  The most frequent targeted shell build
targets were, in order, \texttt{verif\_garbage\_collect}, \texttt{GCGraph},
\texttt{verif\_do\_generation}, and \texttt{gc\_correct}, with 208, 163, 122,
and 119 calls respectively.

The final state has three notable properties.

\begin{itemize}
\item The C collector calls \texttt{forward\_remset} before root forwarding and
scanning.
\item The tracked proof sources no longer rely on the invalid
\texttt{no\_backward\_edge} invariant.
\item The theorem \texttt{garbage\_collect\_isomorphism} is proved with
remembered-set assumptions but concludes ordinary roots-level graph isomorphism.
\end{itemize}

\section{Observations}

\paragraph{The theorem statement is the artifact to defend.}
In a large formal development, a compiling proof is necessary but not sufficient.
The user must defend the statement being proved.  In this case the main danger
was specification drift: strengthening preconditions until the proof became
easier, or letting remembered-set details leak into a theorem whose purpose was
to state graph preservation for roots.  The most important human review was
therefore review of definitions, assumptions, and theorem statements.

\paragraph{LLMs are strong at local invariant plumbing.}
Once a good invariant is chosen, a large amount of proof work is mechanical but
not trivial.  The proof assistant requires exact lemmas about folds, list
updates, graph validity, generation bounds, heap sizes, and compatibility
predicates.  Codex was useful because it could navigate the existing library and
produce these bridge lemmas incrementally.  The resulting proof still benefited
from cleanup: some helper lemmas were introduced only to cross a local gap and
became unnecessary after the proof architecture stabilized.

\paragraph{Existing libraries are both leverage and constraint.}
\VST{} and \CertiGraph{} made the result possible because they already provided
the semantic model, spatial graph representation, and verification discipline.
They also constrained the shape of the proof.  A small change to a definition
could ripple through many files, and automation failures were often symptoms of
larger mismatches between mathematical and spatial invariants.  Codex was more
effective when asked to work within existing patterns than when asked to invent
new ones.

\paragraph{The assistant needs explicit boundaries.}
The collaboration worked best when the human specified boundaries such as
``do not change definitions silently'' and ``do not add unnecessary spec
preconditions.''  These rules made the assistant's search space smaller and made
its changes easier to review.  They also reduced a common risk in LLM-assisted
proof work: proving a nearby theorem rather than the intended theorem.

\paragraph{Cleanup is part of proof completion.}
The branch did not stop at the first successful top-level theorem.  Subsequent
commits removed unused proof helpers, simplified repeated arguments, and updated
deprecated \Coq{}/\Rocq{} names.  For ordinary software this would be routine
maintenance; for mechanized proofs it is also a way to expose the real proof
dependencies.  Without cleanup, it is harder to tell whether the development is
robust or merely compiled once.

\section{Comparison with Prior Experience Reports}

The closest inspiration for this report is
\citet{paraskevopoulou2026machinegenerated}, which describes using Claude Code
to mechanize the ANF transformation correctness proof in CertiCoq.  That work
used the existing CPS proof as a template and reported approximately 7,800 lines
of \Rocq{} proof developed in about 96 hours.  The human authors did not write
proof scripts directly; they supplied high-level guidance and reviewed the
result.

The present experience is similar in three respects.  First, both tasks are
substantial mechanized proofs in established verified-compiler infrastructure.
Second, both rely on the proof assistant as a machine checker, not on the
language model's self-assessment.  Third, both show that an agent can make
progress when it can search a mature codebase, edit files, and run the checker
many times.

The differences are more instructive.  The ANF report is primarily a
template-adaptation story: a known proof technique for CPS was adapted to a
related transformation.  The mutable-\CertiGC{} work is an invariant-repair and
proof-maintenance story.  There was no single adjacent proof to copy.  The old
proof path was invalid because mutable updates changed the mathematical
assumption.  The central problem was to choose a new invariant and thread it
through existing \VST{}, \CertiGraph{}, and collector-correctness layers without
weakening the final theorem.

This distinction affects how one should evaluate the assistant.  For template
adaptation, productivity can be measured by how quickly the assistant instantiates
an existing architecture.  For invariant repair, the main question is whether
the assistant helps converge on the right architecture while preserving the
semantic contract.  In this case, Codex was more like a tireless proof
maintenance collaborator than an autonomous prover.

\section{Limitations and Validation}

This is a single-case experience report, not a controlled experiment.  The work
was performed on a branch that already contained substantial human design and
proof effort.  The repository history provides dates, commits, and diffs, but it
does not by itself identify which lines were written manually and which were
written with Codex.  A more complete version of this report should analyze the
Codex transcripts, subject to privacy and reproducibility constraints.

The strongest validation is machine checking.  The relevant proof files compile,
and the branch was merged upstream.  The development also underwent cleanup
after the theorem was restored; in particular, tracked source files no longer
contain the obsolete \texttt{no\_backward\_edge} route.  Nevertheless, machine
checking validates only the statements that remain in the code.  Human review is
still required to judge whether those statements are the intended ones.

The report also does not claim that Codex independently discovered the
remembered-set proof architecture.  The human role included understanding the
collector semantics, identifying the invalid invariant, and constraining the
form of acceptable specification changes.  The more defensible claim is that an
agentic assistant substantially helped turn that proof-engineering direction
into a compiling, cleaned-up artifact.

\todo{Before submission, add the exact toolchain versions from the final build
and, if possible, a reproducibility note with the commit hash, build command,
and hardware used for the reported build time.}

\section{Conclusion}

The mutable-\CertiGC{} branch shows a pragmatic mode of AI-assisted formal proof
development.  Codex did not replace the need for proof-engineering judgment, but
it was useful once the human supplied the semantic direction and enforced
boundaries around specifications.  The resulting proof completes the mutable
collector story by replacing an invalid no-backward-edge assumption with a
remembered-set invariant while preserving the desired graph-isomorphism theorem.

The broader lesson is that agentic assistants are already useful for large
machine-checked developments when the collaboration is organized around the
checker and around carefully guarded theorem statements.  They are especially
valuable for local repair, invariant propagation, and cleanup in mature
codebases.  The remaining hard work is deciding what should be true.

\bibliographystyle{ACM-Reference-Format}
\bibliography{references}

\end{document}
